%% ==================================================================
%% TOAN12_C4_BAI11.tex -- Toán 12, Chương IV, Bài 11
%% Cấu trúc: Lí thuyết trọng tâm -> Ví dụ minh hoạ (chừa chỗ trống) -> Bài tập luyện tập
%% ==================================================================
\baihoc{11}{Nguyên hàm}

\lithuyet

\begin{dinhnghia}[Nguyên hàm của một hàm số]
Cho hàm số $f(x)$ xác định trên một khoảng $K$ (hoặc một đoạn, hoặc một nửa khoảng). Hàm số $F(x)$ được gọi là một \textbf{nguyên hàm} của hàm số $f(x)$ trên $K$ nếu $F'(x)=f(x)$ với mọi $x$ thuộc $K$.
\end{dinhnghia}

\begin{chuy}
Trường hợp $K=[a;b]$ thì các đẳng thức $F'(a)=f(a)$ và $F'(b)=f(b)$ được hiểu là đạo hàm bên phải tại điểm $x=a$ và đạo hàm bên trái tại điểm $x=b$ của hàm số $F(x)$, tức là
\[\lim_{x\to a^+}\frac{F(x)-F(a)}{x-a}=f(a)\quad\text{và}\quad\lim_{x\to b^-}\frac{F(x)-F(b)}{x-b}=f(b).\]
\end{chuy}

\begin{dinhli}[Họ các nguyên hàm]
Giả sử hàm số $F(x)$ là một nguyên hàm của $f(x)$ trên $K$. Khi đó:
\begin{enumerate}[label=\alph*)]
  \item Với mỗi hằng số $C$, hàm số $F(x)+C$ cũng là một nguyên hàm của $f(x)$ trên $K$;
  \item Nếu hàm số $G(x)$ là một nguyên hàm của $f(x)$ trên $K$ thì tồn tại một hằng số $C$ sao cho $G(x)=F(x)+C$ với mọi $x\in K$.
\end{enumerate}
Như vậy, nếu $F(x)$ là một nguyên hàm của $f(x)$ trên $K$ thì mọi nguyên hàm của $f(x)$ trên $K$ đều có dạng $F(x)+C$ ($C$ là hằng số). Ta gọi $F(x)+C$ ($C\in\mathbb R$) là \textbf{họ các nguyên hàm} của $f(x)$ trên $K$, kí hiệu bởi $\displaystyle\int f(x)\,\mathrm dx$.
\end{dinhli}

\begin{chuy}
\begin{minipage}[c]{.64\linewidth}
\begin{enumerate}[label=\alph*)]
  \item Để tìm họ các nguyên hàm (gọi tắt là tìm nguyên hàm) của hàm số $f(x)$ trên $K$, ta chỉ cần tìm một nguyên hàm $F(x)$ của $f(x)$ trên $K$ và khi đó
  \[\int f(x)\,\mathrm dx=F(x)+C,\ C\text{ là hằng số}.\]
  \item Người ta chứng minh được rằng, nếu hàm số $f(x)$ liên tục trên khoảng $K$ thì $f(x)$ có nguyên hàm trên khoảng đó.
  \item Biểu thức $f(x)\,\mathrm dx$ gọi là vi phân của nguyên hàm $F(x)$, kí hiệu là $\mathrm dF(x)$. Vậy $\mathrm dF(x)=F'(x)\,\mathrm dx=f(x)\,\mathrm dx$.
  \item Khi tìm nguyên hàm của một hàm số mà không chỉ rõ tập $K$, ta hiểu là tìm nguyên hàm của hàm số đó trên tập xác định của nó.
\end{enumerate}
\end{minipage}\hfill
\begin{minipage}[c]{.33\linewidth}\centering
\begin{tikzpicture}[font=\small]
  \node[draw=dcnhan,line width=1pt,regular polygon,regular polygon sides=5,fill=dcvang,inner sep=1pt] (F) at (0,0){$F(x)$};
  \node[draw=dcchinh,line width=1pt,ellipse,fill=dcnen] (f) at (3,0){$f(x)$};
  \draw[-{Stealth},dcnhan,line width=1pt] (F.north) to[bend left=50] node[above,font=\scriptsize,align=center]{Tính\\ đạo hàm} (f.north);
  \draw[-{Stealth},dcchinh,line width=1pt] (f.south) to[bend left=50] node[below,font=\scriptsize,align=center]{Tìm\\ nguyên hàm} (F.south);
\end{tikzpicture}\\[2pt]
{\scriptsize\itshape\color{dcxam}Mối quan hệ giữa đạo hàm và nguyên hàm}
\end{minipage}
\end{chuy}

\begin{congthuc}[Tính chất cơ bản của nguyên hàm]
\begin{itemize}
  \item $\displaystyle\int kf(x)\,\mathrm dx=k\int f(x)\,\mathrm dx\quad(k\ne0)$.
  \item $\displaystyle\int\big[f(x)+g(x)\big]\mathrm dx=\int f(x)\,\mathrm dx+\int g(x)\,\mathrm dx$.
  \item $\displaystyle\int\big[f(x)-g(x)\big]\mathrm dx=\int f(x)\,\mathrm dx-\int g(x)\,\mathrm dx$.
\end{itemize}
\end{congthuc}

\begin{kienthuc}[Hàm số luỹ thừa]
Hàm số $y=x^\alpha$, với $\alpha\in\mathbb R$, được gọi là \emph{hàm số luỹ thừa}. Tập xác định của hàm số luỹ thừa $y=x^\alpha$ tuỳ thuộc vào giá trị của $\alpha$: với $\alpha$ nguyên dương, tập xác định là $\mathbb R$; với $\alpha$ nguyên âm hoặc bằng $0$, tập xác định là $\mathbb R\setminus\{0\}$; với $\alpha$ không nguyên, tập xác định là $(0;+\infty)$.

Hàm số luỹ thừa $y=x^\alpha$ ($\alpha\in\mathbb R$) có đạo hàm với mọi $x>0$ và $\left(x^\alpha\right)'=\alpha x^{\alpha-1}$.
\end{kienthuc}

\begin{congthuc}[Bảng nguyên hàm của một số hàm số thường gặp]
\begin{center}{\renewcommand\arraystretch{2.1}\arrayrulecolor{dcnhan}
\begin{tabular}{|>{\centering\arraybackslash}p{.44\linewidth}|>{\centering\arraybackslash}p{.44\linewidth}|}\hline
$\displaystyle\int0\,\mathrm dx=C$ & $\displaystyle\int1\,\mathrm dx=x+C$\\\hline
$\displaystyle\int x^\alpha\,\mathrm dx=\frac{x^{\alpha+1}}{\alpha+1}+C\ (\alpha\ne-1)$ & $\displaystyle\int\frac1x\,\mathrm dx=\ln|x|+C$\\\hline
$\displaystyle\int\mathrm e^x\,\mathrm dx=\mathrm e^x+C$ & $\displaystyle\int a^x\,\mathrm dx=\frac{a^x}{\ln a}+C\ (0<a\ne1)$\\\hline
$\displaystyle\int\cos x\,\mathrm dx=\sin x+C$ & $\displaystyle\int\sin x\,\mathrm dx=-\cos x+C$\\\hline
$\displaystyle\int\frac{1}{\cos^2x}\,\mathrm dx=\tan x+C$ & $\displaystyle\int\frac{1}{\sin^2x}\,\mathrm dx=-\cot x+C$\\\hline
\end{tabular}}\end{center}
Dựa vào bảng nguyên hàm của các hàm số thường gặp và tính chất cơ bản của nguyên hàm, ta có thể tìm được nguyên hàm của nhiều hàm số khác.
\end{congthuc}

\begin{nhanxet}[Ý nghĩa thực tiễn]
Nếu $S(t)$ là quãng đường vật đi được thì $S'(t)=v(t)$, nên $S(t)$ là một nguyên hàm của hàm vận tốc $v(t)$. Tương tự, doanh thu $R(x)$ là một nguyên hàm của tốc độ biến động doanh thu $M_R(x)=R'(x)$.
\end{nhanxet}

%% ------------------------------------------------ VÍ DỤ
\viduminhhoa
\subsubsection{Tìm nguyên hàm bằng bảng nguyên hàm và tính chất}
\begin{vidu}
Tìm: a) $\displaystyle\int\left(4x^3-3x^2+2x-5\right)\mathrm dx$;\qquad b) $\displaystyle\int\left(\sqrt x+\frac1{x^2}\right)\mathrm dx$ $(x>0)$.
\begin{loigiai}
a) $\displaystyle\int\left(4x^3-3x^2+2x-5\right)\mathrm dx=4\cdot\frac{x^4}4-3\cdot\frac{x^3}3+2\cdot\frac{x^2}2-5x+C=x^4-x^3+x^2-5x+C$.

b) $\displaystyle\int\left(x^{\frac12}+x^{-2}\right)\mathrm dx=\frac{x^{\frac32}}{\frac32}+\frac{x^{-1}}{-1}+C=\frac23x\sqrt x-\frac1x+C$.
\end{loigiai}
\end{vidu}

\begin{vidu}
Tìm: a) $\displaystyle\int\left(3\cos x-\frac{2}{\sin^2x}\right)\mathrm dx$;\qquad b) $\displaystyle\int\left(2\mathrm e^x-3^x\right)\mathrm dx$;\qquad c) $\displaystyle\int\frac{1}{5^x}\,\mathrm dx$.
\begin{loigiai}
a) $\displaystyle\int\left(3\cos x-\frac{2}{\sin^2x}\right)\mathrm dx=3\sin x+2\cot x+C$.

b) $\displaystyle\int\left(2\mathrm e^x-3^x\right)\mathrm dx=2\mathrm e^x-\frac{3^x}{\ln3}+C$.

c) $\displaystyle\int\left(\frac15\right)^x\mathrm dx=\frac{\left(\frac15\right)^x}{\ln\frac15}+C=-\frac{1}{5^x\ln5}+C$.
\end{loigiai}
\end{vidu}

\subsubsection{Biến đổi hàm số về tổng, hiệu các hàm số trong bảng nguyên hàm}
\begin{vidu}
Tìm: a) $\displaystyle\int(2x-1)^2\,\mathrm dx$;\qquad b) $\displaystyle\int\frac{x^2-2x+3}{x}\,\mathrm dx$ $(x>0)$.
\begin{loigiai}
a) $\displaystyle\int\left(4x^2-4x+1\right)\mathrm dx=\frac43x^3-2x^2+x+C$.

b) $\displaystyle\int\left(x-2+\frac3x\right)\mathrm dx=\frac{x^2}2-2x+3\ln x+C$.
\end{loigiai}
\end{vidu}

\begin{vidu}
Tìm: a) $\displaystyle\int4\sin^2\frac x2\,\mathrm dx$;\qquad b) $\displaystyle\int\tan^2x\,\mathrm dx$.
\begin{loigiai}
a) Vì $4\sin^2\dfrac x2=2(1-\cos x)$ nên $\displaystyle\int4\sin^2\frac x2\,\mathrm dx=\int(2-2\cos x)\,\mathrm dx=2x-2\sin x+C$.

b) Vì $\tan^2x=\dfrac1{\cos^2x}-1$ nên $\displaystyle\int\tan^2x\,\mathrm dx=\tan x-x+C$.
\end{loigiai}
\end{vidu}

\subsubsection{Tìm nguyên hàm thoả mãn điều kiện cho trước}
\begin{vidu}
Tìm nguyên hàm $F(x)$ của hàm số $f(x)=3x^2-4x+1$ thoả mãn $F(1)=2$.
\begin{loigiai}
$F(x)=x^3-2x^2+x+C$. Vì $F(1)=1-2+1+C=C=2$ nên $F(x)=x^3-2x^2+x+2$.
\end{loigiai}
\end{vidu}

\begin{vidu}
Cho hàm số $y=f(x)$ xác định trên $(0;+\infty)$, có $f'(x)=2x+\dfrac1{x^2}$ với mọi $x\in(0;+\infty)$ và $f(1)=1$. Tính $f(4)$.
\begin{loigiai}
$f(x)=\displaystyle\int\left(2x+\frac1{x^2}\right)\mathrm dx=x^2-\frac1x+C$. Từ $f(1)=1-1+C=1$ suy ra $C=1$. Vậy $f(4)=16-\dfrac14+1=\dfrac{67}4$.
\end{loigiai}
\end{vidu}

\subsubsection{Bài toán thực tiễn: tìm đại lượng khi biết tốc độ thay đổi}
\begin{vidu}
Một máy bay chạy đà trên đường băng với vận tốc $v(t)=5+3t$ (m/s), $t$ là thời gian (giây) kể từ khi bắt đầu chạy đà. Sau $30$ giây máy bay rời đường băng. Tính quãng đường máy bay đã di chuyển trên đường băng.
\begin{loigiai}
Gọi $S(t)$ là quãng đường máy bay đi được sau $t$ giây. Ta có $S'(t)=v(t)$ nên $S(t)=\displaystyle\int(5+3t)\,\mathrm dt=5t+\frac32t^2+C$. Vì $S(0)=0$ nên $C=0$. Vậy $S(30)=150+1\,350=1\,500$ (m).
\end{loigiai}
\end{vidu}

\begin{vidu}
Tốc độ biến động của doanh thu khi bán $x$ con chip là $M_R(x)=R'(x)=300-0{,}1x$ (triệu đồng/con chip). Tìm doanh thu của công ty khi bán $1\,000$ con chip, biết $R(0)=0$.
\begin{loigiai}
$R(x)=\displaystyle\int(300-0{,}1x)\,\mathrm dx=300x-0{,}05x^2+C$; $R(0)=0\Rightarrow C=0$. Vậy $R(1\,000)=300\,000-50\,000=250\,000$ (triệu đồng).
\end{loigiai}
\end{vidu}

\begin{viduds}
Cho hàm số $f(x)=2x+\mathrm e^x$. Xét tính đúng sai của các mệnh đề sau:
\begin{dungsai}
  \y $\displaystyle\int f(x)\,\mathrm dx=x^2+\mathrm e^x+C$. \oDS
  \y Hàm số $F(x)=x^2+\mathrm e^x-1$ là nguyên hàm của $f(x)$ thoả mãn $F(0)=0$. \oDS
  \y $\displaystyle\int f(x)\,\mathrm dx=2+\mathrm e^x+C$. \oDS
  \y Nếu $F(x)$ là nguyên hàm của $f(x)$ thoả mãn $F(0)=0$ thì $F(1)=\mathrm e+1$. \oDS
\end{dungsai}
\begin{loigiai}
\tuduy Tìm một nguyên hàm rồi kiểm tra từng mệnh đề; mệnh đề c) nhầm giữa đạo hàm và nguyên hàm.

a) \textbf{Đúng}. b) \textbf{Đúng}: $F(0)=0+1-1=0$. c) \textbf{Sai}: $2+\mathrm e^x$ là đạo hàm của $f(x)$. d) \textbf{Sai}: $F(1)=1+\mathrm e-1=\mathrm e$.
\end{loigiai}
\end{viduds}
\subsubsection{Bài toán thực tế}
\begin{vidu}
Một ô tô đang chạy với vận tốc $16$ m/s thì người lái đạp phanh; từ thời điểm đó ô tô chuyển động chậm dần đều với vận tốc $v(t)=-8t+16$ (m/s), $t$ là thời gian (giây) tính từ lúc đạp phanh. Từ lúc đạp phanh đến khi dừng hẳn, ô tô di chuyển được bao nhiêu mét?
\hinh{\begin{tikzpicture}[scale=0.85,font=\footnotesize,line join=round,line cap=round,>=Stealth]
\draw[thick] (-.3,0)--(6.3,0); \fill[gray!25] (-.3,-.25) rectangle (6.3,0);
\draw[dashed,white,thick] (0,-.12)--(6.2,-.12);
\begin{scope}[shift={(0,0)}]\filldraw[fill=gray!35] (0,.18) rectangle (1.3,.55); \filldraw[fill=gray!15] (.25,.55)--(.4,.85)--(.95,.85)--(1.1,.55)--cycle;
\filldraw[fill=black] (.3,.18) circle (.15) (1,.18) circle (.15);\end{scope}
\draw[->,thick] (1.5,.4)--(2.5,.4) node[above,midway]{$v$};
\draw[<->] (0,1.2)--(4.6,1.2) node[midway,above]{quãng đường phanh $s$};
\draw[dashed] (4.6,0)--(4.6,1.3);
\end{tikzpicture}}

\begin{loigiai}
\giai Xe dừng khi $v(t)=0\Leftrightarrow t=2$. $S(t)=\displaystyle\int v(t)\,\mathrm{d}t=-4t^2+16t$ (vì $S(0)=0$), $S(2)=16$ (m). Đáp số: $16$.
\end{loigiai}
\end{vidu}
\begin{vidu}
Một viên đạn được bắn thẳng đứng lên trên từ mặt đất với vận tốc $v(t)=160-9{,}8t$ (m/s). Tính độ cao lớn nhất (mét) mà viên đạn đạt được (làm tròn kết quả đến hàng đơn vị).
\hinh{\begin{tikzpicture}[scale=0.9,font=\footnotesize,line join=round,line cap=round,>=Stealth]
\draw[thick] (-1.6,0)--(1.6,0); \fill[pattern=north east lines] (-1.6,-.15) rectangle (1.6,0);
\draw[dashed] (0,0)--(0,2.8); \fill (0,2.9) ellipse (.05 and .12);
\draw[->,thick] (.25,.4)--(.25,1.6) node[midway,right]{$v(t)$};
\node at (0,3.3){độ cao lớn nhất?}; \draw[->] (-1.1,0)--(-1.1,3) node[above]{$h$};
\end{tikzpicture}}

\begin{loigiai}
\giai $h(t)=160t-4{,}9t^2$; độ cao lớn nhất tại $t=\dfrac{800}{49}$: $h=\dfrac{64000}{49}\approx 1306$. Đáp số: $1306$.
\end{loigiai}
\end{vidu}
\begin{vidu}
Một quần thể vi khuẩn ban đầu có $500$ con. Tốc độ tăng trưởng của quần thể tại thời điểm $t$ (giờ) là $P'(t)=50+2t$ (con/giờ). Số vi khuẩn sau $12$ giờ. Tính kết quả (đơn vị: con).
\begin{loigiai}
\giai $P(t)=500+50t+t^2$ (vì $P(0)=500$). $P(12)=1244$ (con). Đáp số: $1244$.
\end{loigiai}
\end{vidu}
\begin{viduds}
Một viên đạn được bắn thẳng đứng lên trên từ mặt đất với vận tốc $v(t)=49-9{,}8t$ (m/s), $t$ tính bằng giây ($t=0$ là thời điểm bắn). Xét tính đúng sai của các mệnh đề sau:
\hinh{\begin{tikzpicture}[scale=0.9,font=\footnotesize,line join=round,line cap=round,>=Stealth]
\draw[thick] (-1.6,0)--(1.6,0); \fill[pattern=north east lines] (-1.6,-.15) rectangle (1.6,0);
\draw[dashed] (0,0)--(0,2.8); \fill (0,2.9) ellipse (.05 and .12);
\draw[->,thick] (.25,.4)--(.25,1.6) node[midway,right]{$v(t)$};
\node at (0,3.3){độ cao lớn nhất?}; \draw[->] (-1.1,0)--(-1.1,3) node[above]{$h$};
\end{tikzpicture}}

\begin{dungsai}
  \y Viên đạn rơi trở lại mặt đất sau $5$ giây kể từ lúc bắn. \oDS
  \y Độ cao của viên đạn tại thời điểm $t$ là $h(t)=49t-4{,}9t^2$ (m). \oDS
  \y Độ cao của viên đạn tại thời điểm $t$ là $h(t)=49t-9{,}8t^2$ (m). \oDS
  \y Độ cao lớn nhất của viên đạn là $245$ m. \oDS
\end{dungsai}
\begin{loigiai}
\giai $h(t)=\displaystyle\int v(t)\,\mathrm{d}t=49t-4{,}9t^2$ (vì $h(0)=0$); $v(t)=0\Leftrightarrow t=5$. a) Sai: thời gian đúng là $10$ giây. b) Đúng. c) Sai: hệ số đúng là $4{,}9$. d) Sai: giá trị đúng là $122{,}5$ m.
\end{loigiai}
\end{viduds}

%% ------------------------------------------------ BÀI TẬP
\luyentap
\phanbt{Trắc nghiệm nhiều phương án lựa chọn}
\begin{bai}Họ tất cả các nguyên hàm của hàm số $f(x)=3^x$ là
\pa{$x\cdot 3^{x-1}+C$}{$\dfrac{3^{x+1}}{x+1}+C$}{$3^x\ln 3+C$}{$\dfrac{3^x}{\ln 3}+C$}\dapso{D. $\int a^x\,\mathrm dx=\dfrac{a^x}{\ln a}+C$.}\end{bai}
\begin{bai}Họ tất cả các nguyên hàm của hàm số $f(x)=5\cos x-5\sin x$ là
\pa{$-5\sin x-5\cos x+C$}{$5\sin x-5\cos x+C$}{$-5\sin x+5\cos x+C$}{$5\sin x+5\cos x+C$}\dapso{D. $\int\cos x\,\mathrm dx=\sin x+C$, $\int\sin x\,\mathrm dx=-\cos x+C$.}\end{bai}
\begin{bai}Họ tất cả các nguyên hàm của hàm số $f(x)=9 x^{2} + 6 x + 2$ là
\pa{$18 x + 6+C$}{$3 x^{3} + 3 x^{2}+C$}{$3 x^{3} + 3 x^{2} + 2 x+C$}{$9 x^{3} + 6 x^{2} + 2 x+C$}\dapso{C. $\int(9 x^{2} + 6 x + 2)\,\mathrm dx=3 x^{3} + 3 x^{2} + 2 x+C$.}\end{bai}
\begin{bai}Hàm số nào dưới đây là một nguyên hàm của hàm số $f(x)=3 x^{2} + 2$?
\pa{$F(x)=x^{3} + 2 x + 2$}{$F(x)=3 x^{3} + 2 x$}{$F(x)=x^{3} + 2$}{$F(x)=6 x$}\dapso{A. Ta kiểm tra $F'(x)=f(x)$: $\left(x^{3} + 2 x + 2\right)'=3 x^{2} + 2$.}\end{bai}
\begin{bai}Họ tất cả các nguyên hàm của hàm số $f(x)=\dfrac{x^2+4x-1}{x}$ trên khoảng $(0;+\infty)$ là
\pa{$\dfrac{x^2}{2}+4x-\ln|x|+C$}{$x+4-\ln|x|+C$}{$\dfrac{x^2}{2}+4x+\dfrac{1}{x^2}+C$}{$x^2+4x-\ln|x|+C$}\dapso{A. $f(x)=x+4-\cdot\dfrac1x$ nên $\int f(x)\,\mathrm dx=\dfrac{x^2}{2}+4x-\ln|x|+C$.}\end{bai}
\begin{bai}Họ tất cả các nguyên hàm của hàm số $f(x)=\dfrac{1}{\cos^2x}+2$ trên khoảng $\left(-\dfrac\pi2;\dfrac\pi2\right)$ là
\pa{$\tan x+2x+C$}{$\cot x+2x+C$}{$\tan x+2+C$}{$-\tan x+2x+C$}\dapso{A. $\int\dfrac{1}{\cos^2x}\,\mathrm dx=\tan x+C$.}\end{bai}
\begin{bai}Biết $F(x)$ là một nguyên hàm của hàm số $f(x)=6 x + 2$ và $F(0)=1$. Giá trị của $F(2)$ bằng
\pa{$19$}{$13$}{$16$}{$17$}\dapso{D. $F(x)=3 x^{2} + 2 x+C$, $F(0)=1\Rightarrow C=1$. Vậy $F(2)=17$.}\end{bai}
\begin{bai}Họ tất cả các nguyên hàm của hàm số $f(x)=(x+2)(x-1)$ là
\pa{$\dfrac{x^{3}}{3} + \dfrac{x^{2}}{2} + 2 x+C$}{$\left(\dfrac{x^{2}}{2} - x\right) \left(\dfrac{x^{2}}{2} + 2 x\right)+C$}{$x^{3} + x^{2} - 2 x+C$}{$\dfrac{x^{3}}{3} + \dfrac{x^{2}}{2} - 2 x+C$}\dapso{D. $f(x)=x^{2} + x - 2$ nên $\int f(x)\,\mathrm dx=\dfrac{x^{3}}{3} + \dfrac{x^{2}}{2} - 2 x+C$.}\end{bai}
\begin{bai}Một vật chuyển động với vận tốc $v(t)=4 t + 3$ (m/s), trong đó $t$ là thời gian tính bằng giây. Gọi $s(t)$ là quãng đường vật đi được sau $t$ giây, $s(0)=0$. Quãng đường vật đi được sau 5 giây là
\pa{$115$ m}{$65$ m}{$68$ m}{$23$ m}\dapso{B. $s(t)=\int v(t)\,\mathrm dt=2 t^{2} + 3 t+C$, $s(0)=0\Rightarrow C=0$; $s(5)=65$.}\end{bai}
\begin{bai}Một vật chuyển động với vận tốc $v(t)=6 t + 10$ (m/s), trong đó $t$ là thời gian tính bằng giây. Gọi $s(t)$ là quãng đường vật đi được sau $t$ giây, $s(0)=0$. Quãng đường vật đi được sau 2 giây là
\pa{$42$ m}{$22$ m}{$44$ m}{$32$ m}\dapso{D. $s(t)=\int v(t)\,\mathrm dt=3 t^{2} + 10 t+C$, $s(0)=0\Rightarrow C=0$; $s(2)=32$.}\end{bai}
\begin{bai}Nước được bơm vào một bể với tốc độ $r(t)=4t+20$ (lít/phút), $t$ tính bằng phút. Lúc đầu bể không có nước. Sau $15$ phút, lượng nước trong bể là
\hinh{\begin{tikzpicture}[scale=0.9,font=\footnotesize,line join=round,line cap=round,>=Stealth]
\draw[thick] (0,2.2)--(0,0)--(2.6,0)--(2.6,2.2); \fill[gray!25] (0,0) rectangle (2.6,1.1); \draw (0,1.1)--(2.6,1.1);
\draw[thick] (-.4,2.6)--(1,2.6)--(1,2.35); \draw[gray,dashed] (1,2.3)--(1,1.15);
\draw[->] (1.25,2.2)--(1.25,1.6) node[right,midway]{$r(t)$}; \draw[<->] (2.85,0)--(2.85,1.1) node[midway,right]{$V(t)$};
\end{tikzpicture}}

\pa{$450$ lít}{$80$ lít}{$1200$ lít}{$750$ lít}
\dapso{D. Lượng nước $V(t)=\displaystyle\int r(t)\,\mathrm{d}t=2t^2+20t$ (vì $V(0)=0$); $V(15)=750$ (lít).}\end{bai}
\begin{bai}Một viên đạn được bắn thẳng đứng lên trên từ mặt đất với vận tốc $v(t)=120-9{,}8t$ (m/s). Độ cao lớn nhất (mét) mà viên đạn đạt được (làm tròn đến hàng đơn vị) là
\hinh{\begin{tikzpicture}[scale=0.9,font=\footnotesize,line join=round,line cap=round,>=Stealth]
\draw[thick] (-1.6,0)--(1.6,0); \fill[pattern=north east lines] (-1.6,-.15) rectangle (1.6,0);
\draw[dashed] (0,0)--(0,2.8); \fill (0,2.9) ellipse (.05 and .12);
\draw[->,thick] (.25,.4)--(.25,1.6) node[midway,right]{$v(t)$};
\node at (0,3.3){độ cao lớn nhất?}; \draw[->] (-1.1,0)--(-1.1,3) node[above]{$h$};
\end{tikzpicture}}

\pa{$367$ m}{$784$ m}{$735$ m}{$1469$ m}
\dapso{C. $h(t)=120t-4{,}9t^2$; $v(t)=0\Leftrightarrow t=\dfrac{600}{49}$; $h_{\max}=h(\dfrac{600}{49})=\dfrac{36000}{49}\approx 735$ (m).}\end{bai}
\begin{bai}Một viên đạn được bắn thẳng đứng lên trên từ mặt đất với vận tốc $v(t)=49-9{,}8t$ (m/s). Độ cao lớn nhất (mét) mà viên đạn đạt được (làm tròn đến hàng đơn vị) là
\hinh{\begin{tikzpicture}[scale=0.9,font=\footnotesize,line join=round,line cap=round,>=Stealth]
\draw[thick] (-1.6,0)--(1.6,0); \fill[pattern=north east lines] (-1.6,-.15) rectangle (1.6,0);
\draw[dashed] (0,0)--(0,2.8); \fill (0,2.9) ellipse (.05 and .12);
\draw[->,thick] (.25,.4)--(.25,1.6) node[midway,right]{$v(t)$};
\node at (0,3.3){độ cao lớn nhất?}; \draw[->] (-1.1,0)--(-1.1,3) node[above]{$h$};
\end{tikzpicture}}

\pa{$172$ m}{$61$ m}{$245$ m}{$123$ m}
\dapso{D. $h(t)=49t-4{,}9t^2$; $v(t)=0\Leftrightarrow t=5$; $h_{\max}=h(5)=\dfrac{245}{2}\approx 123$ (m).}\end{bai}
\begin{bai}Một viên đạn được bắn thẳng đứng lên trên từ mặt đất với vận tốc $v(t)=196-9{,}8t$ (m/s). Độ cao lớn nhất (mét) mà viên đạn đạt được (làm tròn đến hàng đơn vị) là
\hinh{\begin{tikzpicture}[scale=0.9,font=\footnotesize,line join=round,line cap=round,>=Stealth]
\draw[thick] (-1.6,0)--(1.6,0); \fill[pattern=north east lines] (-1.6,-.15) rectangle (1.6,0);
\draw[dashed] (0,0)--(0,2.8); \fill (0,2.9) ellipse (.05 and .12);
\draw[->,thick] (.25,.4)--(.25,1.6) node[midway,right]{$v(t)$};
\node at (0,3.3){độ cao lớn nhất?}; \draw[->] (-1.1,0)--(-1.1,3) node[above]{$h$};
\end{tikzpicture}}

\pa{$1960$ m}{$3920$ m}{$980$ m}{$2009$ m}
\dapso{A. $h(t)=196t-4{,}9t^2$; $v(t)=0\Leftrightarrow t=20$; $h_{\max}=h(20)=1960\approx 1960$ (m).}\end{bai}

\phanbt{Trắc nghiệm đúng/sai}
\begin{bai}Cho hàm số $f(x)=6 x^{2} - 4 x + 1$. Xét tính đúng sai của các mệnh đề sau:
\begin{dungsai}
  \y $\int f(x)\,\mathrm dx=6 x^{3} - 6 x^{2} + 3 x+C$. \oDS
  \y Hàm số $F(x)=2 x^{3} - 2 x^{2} + x+2$ là một nguyên hàm của $f(x)$. \oDS
  \y $\int f'(x)\,\mathrm dx=12 x - 4+C$. \oDS
  \y Nếu $F(x)$ là nguyên hàm của $f(x)$ thoả $F(0)=2$ thì $F(2)=14$. \oDS
\end{dungsai}
\dapso{a) S. b) Đ. c) S. d) S. $\int f(x)\,\mathrm dx=2 x^{3} - 2 x^{2} + x+C$; $\int f'(x)\,\mathrm dx=f(x)+C$; $F(x)=2 x^{3} - 2 x^{2} + x + 2$, $F(2)=12$.}\end{bai}
\begin{bai}Một vật chuyển động thẳng với vận tốc $v(t)=3 t^{2} + 2 t$ (m/s), $t$ là thời gian tính bằng giây, quãng đường tính từ thời điểm $t=0$.
\begin{dungsai}
  \y Quãng đường vật đi được sau $t$ giây là $s(t)=t^{3} + t^{2}$ (m). \oDS
  \y Gia tốc của vật tại thời điểm $t$ là $a(t)=6 t + 2$ (m/s$^2$). \oDS
  \y Quãng đường vật đi được trong 2 giây đầu tiên là $14$ m. \oDS
  \y Gia tốc của vật tại thời điểm $t=1$ giây là $8$ m/s$^2$. \oDS
\end{dungsai}
\dapso{a) Đ. b) Đ. c) S. d) Đ. $s(t)=\int v(t)\,\mathrm dt=t^{3} + t^{2}$ (vì $s(0)=0$); $a(t)=v'(t)=6 t + 2$; $s(2)=12$.}\end{bai}
\begin{bai}Một viên đạn được bắn thẳng đứng lên trên từ mặt đất với vận tốc $v(t)=98-9{,}8t$ (m/s), $t$ tính bằng giây ($t=0$ là thời điểm bắn). Xét tính đúng sai của các mệnh đề sau:
\hinh{\begin{tikzpicture}[scale=0.9,font=\footnotesize,line join=round,line cap=round,>=Stealth]
\draw[thick] (-1.6,0)--(1.6,0); \fill[pattern=north east lines] (-1.6,-.15) rectangle (1.6,0);
\draw[dashed] (0,0)--(0,2.8); \fill (0,2.9) ellipse (.05 and .12);
\draw[->,thick] (.25,.4)--(.25,1.6) node[midway,right]{$v(t)$};
\node at (0,3.3){độ cao lớn nhất?}; \draw[->] (-1.1,0)--(-1.1,3) node[above]{$h$};
\end{tikzpicture}}

\begin{dungsai}
  \y Độ cao lớn nhất của viên đạn là $490$ m. \oDS
  \y Viên đạn rơi trở lại mặt đất sau $10$ giây kể từ lúc bắn. \oDS
  \y Độ cao của viên đạn tại thời điểm $t$ là $h(t)=98t-9{,}8t^2$ (m). \oDS
  \y Độ cao của viên đạn tại thời điểm $t$ là $h(t)=98t-4{,}9t^2$ (m). \oDS
\end{dungsai}
\dapso{$h(t)=\displaystyle\int v(t)\,\mathrm{d}t=98t-4{,}9t^2$ (vì $h(0)=0$); $v(t)=0\Leftrightarrow t=10$. a) Đúng. b) Sai: thời gian đúng là $20$ giây. c) Sai: hệ số đúng là $4{,}9$. d) Đúng.}\end{bai}
\begin{bai}Một viên đạn được bắn thẳng đứng lên trên từ mặt đất với vận tốc $v(t)=245-9{,}8t$ (m/s), $t$ tính bằng giây ($t=0$ là thời điểm bắn). Xét tính đúng sai của các mệnh đề sau:
\hinh{\begin{tikzpicture}[scale=0.9,font=\footnotesize,line join=round,line cap=round,>=Stealth]
\draw[thick] (-1.6,0)--(1.6,0); \fill[pattern=north east lines] (-1.6,-.15) rectangle (1.6,0);
\draw[dashed] (0,0)--(0,2.8); \fill (0,2.9) ellipse (.05 and .12);
\draw[->,thick] (.25,.4)--(.25,1.6) node[midway,right]{$v(t)$};
\node at (0,3.3){độ cao lớn nhất?}; \draw[->] (-1.1,0)--(-1.1,3) node[above]{$h$};
\end{tikzpicture}}

\begin{dungsai}
  \y Độ cao của viên đạn tại thời điểm $t$ là $h(t)=245t-9{,}8t^2$ (m). \oDS
  \y Viên đạn rơi trở lại mặt đất sau $25$ giây kể từ lúc bắn. \oDS
  \y Độ cao lớn nhất của viên đạn là $6125$ m. \oDS
  \y Độ cao của viên đạn tại thời điểm $t$ là $h(t)=245t-4{,}9t^2$ (m). \oDS
\end{dungsai}
\dapso{$h(t)=\displaystyle\int v(t)\,\mathrm{d}t=245t-4{,}9t^2$ (vì $h(0)=0$); $v(t)=0\Leftrightarrow t=25$. a) Sai: hệ số đúng là $4{,}9$. b) Sai: thời gian đúng là $50$ giây. c) Sai: giá trị đúng là $3062{,}5$ m. d) Đúng.}\end{bai}

\phanbt{Trả lời ngắn}
\begin{bai}Biết $F(x)$ là một nguyên hàm của hàm số $f(x)=3 x^{2} - 2 x$ và $F(0)=3$. Tính $F(2)$.\otraloi\dapso{7. $F(x)=x^{3} - x^{2} + 3$, $F(2)=7$.}\end{bai}
\begin{bai}Biết $F(x)$ là một nguyên hàm của hàm số $f(x)=3 x^{2} + 4 x$ và $F(0)=5$. Tính $F(3)$.\otraloi\dapso{50. $F(x)=x^{3} + 2 x^{2} + 5$, $F(3)=50$.}\end{bai}
\begin{bai}Biết $F(x)$ là một nguyên hàm của hàm số $f(x)=6 x^{2} + 2 x$ và $F(0)=1$. Tính $F(2)$.\otraloi\dapso{21. $F(x)=2 x^{3} + x^{2} + 1$, $F(2)=21$.}\end{bai}
\begin{bai}Biết $F(x)$ là một nguyên hàm của hàm số $f(x)=\mathrm e^x+2x$ và $F(0)=1$. Tính $F(1)$ (làm tròn đến hàng phần trăm).\otraloi\dapso{3{,}72. $F(x)=\mathrm e^x+x^2+C$, $F(0)=1+C=1\Rightarrow C=0$; $F(1)=\mathrm e+1\approx3{,}72$.}\end{bai}
\begin{bai}Một vật chuyển động có gia tốc $a(t)=6 t + 3$ (m/s$^2$), $t$ tính bằng giây. Biết vận tốc ban đầu của vật là 5 m/s. Tính vận tốc của vật tại thời điểm $t=2$ giây (m/s).\otraloi\dapso{23. $v(t)=3 t^{2} + 3 t+5$; $v(2)=23$.}\end{bai}
\begin{bai}Biết $F(x)$ là một nguyên hàm của hàm số $f(x)=\mathrm e^x+2x$ và $F(0)=3$. Tính $F(1)$ (làm tròn đến hàng phần trăm).\otraloi\dapso{5{,}72. $F(x)=\mathrm e^x+x^2+C$, $F(0)=1+C=3\Rightarrow C=2$; $F(1)=\mathrm e+3\approx5{,}72$.}\end{bai}
\begin{bai}Một ô tô đang chạy với vận tốc $15$ m/s thì người lái đạp phanh; từ thời điểm đó ô tô chuyển động chậm dần đều với vận tốc $v(t)=-3t+15$ (m/s), $t$ là thời gian (giây) tính từ lúc đạp phanh. Từ lúc đạp phanh đến khi dừng hẳn, ô tô di chuyển được bao nhiêu mét?
\hinh{\begin{tikzpicture}[scale=0.85,font=\footnotesize,line join=round,line cap=round,>=Stealth]
\draw[thick] (-.3,0)--(6.3,0); \fill[gray!25] (-.3,-.25) rectangle (6.3,0);
\draw[dashed,white,thick] (0,-.12)--(6.2,-.12);
\begin{scope}[shift={(0,0)}]\filldraw[fill=gray!35] (0,.18) rectangle (1.3,.55); \filldraw[fill=gray!15] (.25,.55)--(.4,.85)--(.95,.85)--(1.1,.55)--cycle;
\filldraw[fill=black] (.3,.18) circle (.15) (1,.18) circle (.15);\end{scope}
\draw[->,thick] (1.5,.4)--(2.5,.4) node[above,midway]{$v$};
\draw[<->] (0,1.2)--(4.6,1.2) node[midway,above]{quãng đường phanh $s$};
\draw[dashed] (4.6,0)--(4.6,1.3);
\end{tikzpicture}}
\otraloi
\dapso{37,5. Xe dừng khi $v(t)=0\Leftrightarrow t=5$. $S(t)=\displaystyle\int v(t)\,\mathrm{d}t=-\dfrac{3}{2}t^2+15t$ (vì $S(0)=0$), $S(5)=\dfrac{75}{2}$ (m).}\end{bai}
\begin{bai}Một viên đạn được bắn thẳng đứng lên trên từ mặt đất với vận tốc $v(t)=100-9{,}8t$ (m/s). Tính độ cao lớn nhất (mét) mà viên đạn đạt được (làm tròn kết quả đến hàng đơn vị).
\hinh{\begin{tikzpicture}[scale=0.9,font=\footnotesize,line join=round,line cap=round,>=Stealth]
\draw[thick] (-1.6,0)--(1.6,0); \fill[pattern=north east lines] (-1.6,-.15) rectangle (1.6,0);
\draw[dashed] (0,0)--(0,2.8); \fill (0,2.9) ellipse (.05 and .12);
\draw[->,thick] (.25,.4)--(.25,1.6) node[midway,right]{$v(t)$};
\node at (0,3.3){độ cao lớn nhất?}; \draw[->] (-1.1,0)--(-1.1,3) node[above]{$h$};
\end{tikzpicture}}
\otraloi
\dapso{510. $h(t)=100t-4{,}9t^2$; độ cao lớn nhất tại $t=\dfrac{500}{49}$: $h=\dfrac{25000}{49}\approx 510$.}\end{bai}
\begin{bai}Một ô tô đang chạy với vận tốc $20$ m/s thì người lái đạp phanh; từ thời điểm đó ô tô chuyển động chậm dần đều với vận tốc $v(t)=-4t+20$ (m/s), $t$ là thời gian (giây) tính từ lúc đạp phanh. Từ lúc đạp phanh đến khi dừng hẳn, ô tô di chuyển được bao nhiêu mét?
\hinh{\begin{tikzpicture}[scale=0.85,font=\footnotesize,line join=round,line cap=round,>=Stealth]
\draw[thick] (-.3,0)--(6.3,0); \fill[gray!25] (-.3,-.25) rectangle (6.3,0);
\draw[dashed,white,thick] (0,-.12)--(6.2,-.12);
\begin{scope}[shift={(0,0)}]\filldraw[fill=gray!35] (0,.18) rectangle (1.3,.55); \filldraw[fill=gray!15] (.25,.55)--(.4,.85)--(.95,.85)--(1.1,.55)--cycle;
\filldraw[fill=black] (.3,.18) circle (.15) (1,.18) circle (.15);\end{scope}
\draw[->,thick] (1.5,.4)--(2.5,.4) node[above,midway]{$v$};
\draw[<->] (0,1.2)--(4.6,1.2) node[midway,above]{quãng đường phanh $s$};
\draw[dashed] (4.6,0)--(4.6,1.3);
\end{tikzpicture}}
\otraloi
\dapso{50. Xe dừng khi $v(t)=0\Leftrightarrow t=5$. $S(t)=\displaystyle\int v(t)\,\mathrm{d}t=-2t^2+20t$ (vì $S(0)=0$), $S(5)=50$ (m).}\end{bai}
\phieutln

\phanbt{Tự luận}
\begin{bai}Tìm: a) $\displaystyle\int\left(3x^2+2x-1\right)\mathrm dx$;\quad b) $\displaystyle\int\left(2x-\frac1x\right)^2\mathrm dx$;\quad c) $\displaystyle\int\left(2^x+\frac3{x^2}\right)\mathrm dx$.
\dongke{6}
\dapso{a) $x^3+x^2-x+C$. b) $\dfrac43x^3-4x-\dfrac1x+C$. c) $\dfrac{2^x}{\ln2}-\dfrac3x+C$.}\end{bai}
\begin{bai}Tìm một nguyên hàm $F(x)$ của hàm số $f(x)=2\cos x+\dfrac1{\sin^2x}$ thoả mãn $F\left(\dfrac\pi4\right)=-1$.
\dongke{4}
\dapso{$F(x)=2\sin x-\cot x+C$; $F\left(\frac\pi4\right)=\sqrt2-1+C=-1\Rightarrow C=-\sqrt2$.}\end{bai}
\begin{bai}Một viên đạn được bắn thẳng đứng lên trên từ mặt đất với vận tốc $v(t)=160-9{,}8t$ (m/s). Tìm độ cao của viên đạn: a) sau $t=5$ giây; b) khi nó đạt độ cao lớn nhất (làm tròn đến hàng phần mười).
\dongke{6}
\dapso{$h(t)=160t-4{,}9t^2$. a) $h(5)=677{,}5$ m. b) $v=0\Leftrightarrow t=\frac{160}{9{,}8}$, $h_{\max}\approx1\,306{,}1$ m.}\end{bai}
\begin{bai}Cá hồi bơi ngược dòng sông với vận tốc $v(t)=-\dfrac{2t}5+4$ (km/h). Coi $t=0$ là lúc cá bắt đầu bơi vào dòng sông. Tính khoảng cách xa nhất mà con cá có thể bơi được.
\dongke{5}
\dapso{Cá dừng khi $v=0\Leftrightarrow t=10$; $s=\displaystyle\int_0^{10}\left(4-\frac{2t}5\right)\mathrm dt$ hay $s(t)=4t-\frac{t^2}5$, $s(10)=20$ km.}\end{bai}
\begin{bai}Nước được bơm vào một bể với tốc độ $r(t)=4t+10$ (lít/phút), $t$ tính bằng phút. Lúc đầu bể không có nước. Sau $10$ phút, lượng nước trong bể. Tính kết quả (đơn vị: lít).
\hinh{\begin{tikzpicture}[scale=0.9,font=\footnotesize,line join=round,line cap=round,>=Stealth]
\draw[thick] (0,2.2)--(0,0)--(2.6,0)--(2.6,2.2); \fill[gray!25] (0,0) rectangle (2.6,1.1); \draw (0,1.1)--(2.6,1.1);
\draw[thick] (-.4,2.6)--(1,2.6)--(1,2.35); \draw[gray,dashed] (1,2.3)--(1,1.15);
\draw[->] (1.25,2.2)--(1.25,1.6) node[right,midway]{$r(t)$}; \draw[<->] (2.85,0)--(2.85,1.1) node[midway,right]{$V(t)$};
\end{tikzpicture}}

\dongke{6}
\dapso{Lượng nước $V(t)=\displaystyle\int r(t)\,\mathrm{d}t=2t^2+10t$ (vì $V(0)=0$); $V(10)=300$ (lít).}\end{bai}
\begin{bai}Một viên đạn được bắn thẳng đứng lên trên từ mặt đất với vận tốc $v(t)=150-9{,}8t$ (m/s). Tính độ cao lớn nhất (mét) mà viên đạn đạt được (làm tròn kết quả đến hàng phần mười).
\hinh{\begin{tikzpicture}[scale=0.9,font=\footnotesize,line join=round,line cap=round,>=Stealth]
\draw[thick] (-1.6,0)--(1.6,0); \fill[pattern=north east lines] (-1.6,-.15) rectangle (1.6,0);
\draw[dashed] (0,0)--(0,2.8); \fill (0,2.9) ellipse (.05 and .12);
\draw[->,thick] (.25,.4)--(.25,1.6) node[midway,right]{$v(t)$};
\node at (0,3.3){độ cao lớn nhất?}; \draw[->] (-1.1,0)--(-1.1,3) node[above]{$h$};
\end{tikzpicture}}

\dongke{6}
\dapso{$h(t)=150t-4{,}9t^2$; độ cao lớn nhất tại $t=\dfrac{750}{49}$: $h=\dfrac{56250}{49}\approx 1148$.}\end{bai}

\phanbt{Bài tập tự luận rèn luyện}
\begin{bai}Tìm nguyên hàm là phép toán ngược của phép lấy đạo hàm. Hãy tìm các họ nguyên hàm sau và tự kiểm tra lại kết quả bằng cách lấy đạo hàm.
\begin{enumerate}[label=\alph*)]
\item $\displaystyle\int(3x^2-4x+5)\,\mathrm dx$.
\item Tìm nguyên hàm $F(x)$ của hàm số $f(x)=3x^2-4x+5$ thoả mãn $F(1)=2$.
\item $\displaystyle\int(2x+1)^2\,\mathrm dx$ (gợi ý: khai triển trước khi lấy nguyên hàm).
\item $\displaystyle\int\dfrac{x^2+1}{x^2}\,\mathrm dx$ trên khoảng $(0;+\infty)$.
\item $\displaystyle\int\left(\sqrt x+\dfrac1{\sqrt x}\right)\mathrm dx$ trên khoảng $(0;+\infty)$.
\item Bạn Hùng cho rằng $\displaystyle\int x(x+1)\,\mathrm dx=\dfrac{x^2}2\cdot\dfrac{(x+1)^2}2+C$. Hãy chỉ ra chỗ sai và sửa lại cho đúng.
\end{enumerate}
\dapso{a) $x^3-2x^2+5x+C$. b) $F(x)=x^3-2x^2+5x-2$. c) $\dfrac43x^3+2x^2+x+C$. d) $x-\dfrac1x+C$. e) $\dfrac23x\sqrt x+2\sqrt x+C$. f) Không có quy tắc “nguyên hàm của tích bằng tích các nguyên hàm”; đúng là $\displaystyle\int(x^2+x)\,\mathrm dx=\dfrac{x^3}3+\dfrac{x^2}2+C$.}\end{bai}

\begin{bai}Tìm các họ nguyên hàm sau. Ở một số câu, em cần biến đổi biểu thức bằng các công thức lượng giác hoặc tính chất của luỹ thừa trước khi áp dụng bảng nguyên hàm.
\begin{enumerate}[label=\alph*)]
\item $\displaystyle\int\left(\mathrm e^x+2^x\right)\mathrm dx$.
\item $\displaystyle\int(3\cos x-2\sin x)\,\mathrm dx$.
\item $\displaystyle\int\left(\dfrac2{\sin^2x}+\dfrac1x\right)\mathrm dx$ trên khoảng $(0;\pi)$.
\item $\displaystyle\int\tan^2x\,\mathrm dx$ trên khoảng $\left(-\dfrac\pi2;\dfrac\pi2\right)$.
\item $\displaystyle\int\sin^2\dfrac x2\,\mathrm dx$.
\item $\displaystyle\int3^{2x}\,\mathrm dx$.
\end{enumerate}
\dapso{a) $\mathrm e^x+\dfrac{2^x}{\ln2}+C$. b) $3\sin x+2\cos x+C$. c) $-2\cot x+\ln x+C$. d) $\tan^2x=\dfrac1{\cos^2x}-1$: $\tan x-x+C$. e) $\sin^2\dfrac x2=\dfrac{1-\cos x}2$: $\dfrac x2-\dfrac{\sin x}2+C$. f) $3^{2x}=9^x$: $\dfrac{9^x}{\ln9}+C$.}\end{bai}

\begin{bai}Trong mỗi trường hợp dưới đây, người ta chỉ cho biết đạo hàm của hàm số và giá trị của nó tại một điểm. Em hãy “khôi phục” lại hàm số ban đầu.
\begin{enumerate}[label=\alph*)]
\item Tìm $f(x)$ biết $f'(x)=6x-2$ và đồ thị hàm số $y=f(x)$ đi qua điểm $(1;4)$.
\item Với hàm số $f$ ở câu a), tính $f(2)$.
\item Tìm $g(x)$ biết $g'(x)=\mathrm e^x-1$ và $g(0)=3$.
\item Tìm $h(x)$ trên $(0;+\infty)$ biết $h'(x)=\dfrac1x$ và $h(\mathrm e)=2$.
\item Tìm $k(x)$ biết $k'(x)=\cos x$ và $k\left(\dfrac\pi2\right)=0$.
\item Hai bạn cùng tìm nguyên hàm của $f'(x)=6x-2$ nhưng một bạn được $3x^2-2x+3$, bạn kia được $3x^2-2x-5$. Đồ thị của hai hàm số này có quan hệ gì với nhau?
\end{enumerate}
\dapso{a) $f(x)=3x^2-2x+3$. b) $11$. c) $g(x)=\mathrm e^x-x+2$. d) $h(x)=\ln x+1$. e) $k(x)=\sin x-1$. f) Hai đồ thị nhận được từ nhau qua phép tịnh tiến theo phương thẳng đứng $8$ đơn vị.}\end{bai}

\begin{bai}Một bể nước ban đầu chứa sẵn $50$ lít. Người ta mở vòi bơm thêm nước vào bể, tốc độ bơm sau $t$ phút là $r(t)=2t+10$ (lít/phút), tức là vòi chảy ngày càng mạnh. Gọi $V(t)$ là lượng nước trong bể sau $t$ phút.
\hinh{\begin{tikzpicture}[scale=0.9,font=\footnotesize,line join=round,line cap=round,>=Stealth]
\draw[thick] (0,2.2)--(0,0)--(2.6,0)--(2.6,2.2); \fill[gray!25] (0,0) rectangle (2.6,1.1); \draw (0,1.1)--(2.6,1.1);
\draw[thick] (-.4,2.6)--(1,2.6)--(1,2.35); \draw[gray,dashed] (1,2.3)--(1,1.15);
\draw[->] (1.25,2.2)--(1.25,1.6) node[right,midway]{$r(t)$}; \draw[<->] (2.85,0)--(2.85,1.1) node[midway,right]{$V(t)$};
\end{tikzpicture}}
\begin{enumerate}[label=\alph*)]
\item Giải thích vì sao $V'(t)=r(t)$, rồi tìm công thức của $V(t)$.
\item Sau $10$ phút trong bể có bao nhiêu lít nước?
\item Trong khoảng từ phút thứ $5$ đến phút thứ $10$, lượng nước bơm vào bể là bao nhiêu?
\item Sau bao lâu thì bể có $500$ lít nước (làm tròn đến hàng phần trăm)?
\item Dung tích của bể là $1\,000$ lít. Sau bao lâu thì bể đầy (làm tròn đến hàng phần trăm)?
\item Tại thời điểm bể vừa đầy, vòi đang chảy với tốc độ bao nhiêu (làm tròn đến hàng phần trăm)?
\end{enumerate}
\dapso{a) $V'(t)$ là tốc độ thay đổi lượng nước; $V(t)=t^2+10t+50$. b) $250$ lít. c) $V(10)-V(5)=125$ lít. d) $t^2+10t-450=0\Rightarrow t=-5+\sqrt{475}\approx16{,}79$ phút. e) $t=-5+\sqrt{975}\approx26{,}22$ phút. f) $r\approx62{,}45$ lít/phút.}\end{bai}

\begin{bai}Một chất điểm bắt đầu chuyển động thẳng từ vị trí gốc với vận tốc ban đầu $2$ m/s. Trong quá trình chuyển động, gia tốc của nó thay đổi theo thời gian: $a(t)=6t-4$ (m/s$^2$), $t$ tính bằng giây.
\begin{enumerate}[label=\alph*)]
\item Tìm vận tốc $v(t)$ của chất điểm, biết $v(0)=2$.
\item Tính vận tốc của chất điểm tại thời điểm $t=2$ giây.
\item Tìm toạ độ $s(t)$ của chất điểm, biết $s(0)=0$.
\item Sau $3$ giây, chất điểm cách vị trí xuất phát bao xa?
\item Tìm vận tốc nhỏ nhất của chất điểm và thời điểm đạt được vận tốc đó.
\item Tính quãng đường chất điểm đi được trong giây thứ ba (từ $t=2$ đến $t=3$).
\end{enumerate}
\dapso{a) $v(t)=3t^2-4t+2$. b) $6$ m/s. c) $s(t)=t^3-2t^2+2t$. d) $15$ m. e) $v'(t)=0\Leftrightarrow t=\dfrac23$; $v_{\min}=\dfrac23$ m/s (vận tốc luôn dương nên vật không đổi chiều). f) $s(3)-s(2)=15-4=11$ m.}\end{bai}
